\section{Architecture}
\label{sec:architecture}

We build on the Diffusion Transformer (DiT)~\cite{peebles2023scalable} and extend it with image-based conditioning, conditioning on the next waypoint $t_\text{next} = t_{i^*+1}$, and an optional Brownian-bridge drift term. The output is the estimated drift from Equation \ref{eqn:simplified_dsm_main}. \textit{All methods compared in the paper use the same architecture.}

\paragraph{Inputs and embeddings.} In addition to the noisy latent $x_t \in \mathbb{R}^{C \times H \times W}$, diffusion timestep $t \in [0,1]$, and class label $y$ used by DiT, our model takes a next-waypoint timestep $t_{\text{next}} \in [0,1]$, a set of $L$ conditioning frames $\{x^{(\ell)}\}_{\ell=1}^{L}$ with associated timesteps $\{t^{(\ell)}\}$, and a boolean validity mask $m \in \{0,1\}^{L}$ (with $m_0 = 1$ always). The conditioning frames are patchified by a dedicated embedder $\mathrm{PatchEmbed}_{\text{cond}}$ and additively tagged with the shared sin-cos positional embedding plus a timestep embedding $\phi_{\text{cond}}(t^{(\ell)})$, yielding tokens $z^{(\ell)} \in \mathbb{R}^{T \times D}$. Timesteps are rescaled by a factor of $1000$ before sinusoidal embedding. The AdaLN conditioning vector becomes
\[
c \;=\; \tfrac{1}{2}\bigl(\phi_t(t) + \phi_{t_{\text{next}}}(t_{\text{next}})\bigr) + \phi_y(y),
\]
where $\phi_t$, $\phi_{t_{\text{next}}}$, and $\phi_{\text{cond}}$ are independent timestep MLPs. %We set $\texttt{learn\_sigma}=\text{False}$ and predict only the mean.

\paragraph{Block structure.} Each block extends the DiT block with a cross-attention sub-layer inserted between self-attention and the MLP:
\begin{align}
&x \leftarrow x + g_{\text{sa}}\,\mathrm{SA}(\mathrm{mod}(x; s_{\text{sa}}, \sigma_{\text{sa}})), \\
&x \leftarrow x + g_{\text{xa}}\,\mathrm{XA}\bigl(\mathrm{mod}(x; s_{\text{xa}}, \sigma_{\text{xa}}),\; \mathrm{mod}(z; s_z, \sigma_z)\bigr), \\
&x \leftarrow x + g_{\text{mlp}}\,\mathrm{MLP}(\mathrm{mod}(x; s_{\text{mlp}}, \sigma_{\text{mlp}})),
\end{align}


so the AdaLN-Zero MLP emits $11D$ modulation parameters per block (vs.\ $6D$ in DiT). Two additional LayerNorms are introduced for the cross-attention query and key/value streams. We apply RMSNorm \cite{zhang2019root} to the per-head queries and keys in both self- and cross-attention to stabilize training on long contexts (the original DiT did not include this, but we found the exclusion to be very unstable).

\paragraph{Variable-length cross-attention.} For cross-attention, we use PyTorch's \texttt{varlen\_attn}, which is based on FlashAttention \cite{dao2022flashattention, dao2023flashattention}. %Because each example attends to a different number of valid conditioning frames, we pack frames across the batch and use variable-length attention (\texttt{varlen\_attn}). Each query frame of $T$ tokens attends to the $T \cdot \sum_\ell m_\ell$ tokens of its valid conditioning frames, with cumulative sequence offsets $\mathrm{cu\_seq}_q$ and $\mathrm{cu\_seq}_{kv}$ computed from $m$.

\paragraph{Brownian-bridge drift.}
We have an option to add the Brownian Bridge drift specified in Section \ref{subsec:any_subset}. We add a small epsilon of $10^{-7}$ to the denominator for stability.

\paragraph{Initialization.} We follow DiT's AdaLN-Zero initialization, additionally applying Xavier initialization \cite{glorot2010understanding} to the conditioning patch embedder and $\mathcal{N}(0, 0.02^2)$ to the new timestep MLPs ($\phi_{\text{cond}}$ and $\phi_{t_{\text{next}}}$).

\paragraph{Model Size}
We use the DiT-B model size. With our cross-attention layers and other architectural modifications, this works out to 202,192,912 parameters.

\paragraph{EMA} We use EMA 0.9999, following DiT defaults.
